\section{Analisis Memori dan Stack serta Metodologi Debugging}
\label{sec:memori-stack-metodologi}

Menelusuri register saja belum cukup. Banyak kesalahan hanya tampak
ketika kita memeriksa isi memori: variabel yang tertimpa, indeks larik
yang melewati batas, atau alamat kembali yang sudah bergeser di dalam
\eng{stack}. Bagian ini membahas cara membaca keluaran \code{D} dan
cara memakainya untuk menemukan tiga kesalahan yang paling sering
terjadi.

\subsection{Membaca Keluaran \code{D}}

Perintah \code{D} menampilkan 128 byte memori dalam delapan baris.
Setiap baris memuat tiga bagian yang dibaca secara berbeda: alamat
dalam bentuk segmen:offset, enam belas pasang digit heksadesimal berisi
isi memori (setiap dua digit adalah satu byte), dan kolom karakter
ASCII di sebelah kanan. Byte yang bukan karakter cetak ditampilkan
sebagai titik, sehingga kolom karakter paling cepat dipakai untuk
mengenali letak sebuah string.

Perhatikan contoh pada Kode~\ref{lst:sesi-debug-lengkap}:

\begin{lstlisting}[numbers=none, backgroundcolor=\color{obeGrayLight}]
13A6:0000 4A 75 6D 6C 61 68 20 31-2E 2E 35 20 3D 20 24 00  Jumlah 1..5 = $.
\end{lstlisting}

Kolom karakter memperlihatkan string \code{Jumlah 1..5 = \$},
sedangkan kolom heksadesimal memperlihatkan kode ASCII-nya: \hex{4A}
untuk \code{J}, \hex{75} untuk \code{u}, dan seterusnya sampai
\hex{24} untuk tanda \code{\$}. Byte terakhir bernilai \hex{00},
yaitu penanda akhir yang dipasang assembler setelah string.

\subsection{Memeriksa Variabel dengan \code{D}}

Program pada Kode~\ref{lst:off-by-one} bermaksud menjumlahkan lima
byte pada larik \code{nilai}, tetapi hasilnya tidak sesuai harapan.

\begin{lstlisting}[
    language=MASM,
    caption={Program dengan kesalahan \eng{off-by-one} pada larik.},
    label={lst:off-by-one}
]
.DATA
    nilai   DB  10, 20, 30, 40, 50      ; lima byte data

.CODE
START:
    MOV AX, @DATA
    MOV DS, AX
    XOR DX, DX                          ; DX = akumulator jumlah
    MOV CX, 6                           ; BUG: seharusnya MOV CX, 5
    MOV SI, 0                           ; SI = indeks larik

ULANG:
    MOV AH, 0                           ; nolkan AH
    MOV AL, nilai[SI]                   ; ambil elemen ke-SI
    ADD DX, AX                          ; tambahkan ke akumulator
    INC SI                              ; naikkan indeks
    LOOP ULANG                          ; ulangi sebanyak CX

    MOV AH, 4Ch
    MOV AL, 00h
    INT 21h
END START
\end{lstlisting}

Jumlah yang benar adalah \dec{150} atau \hex{0096}. Untuk memeriksa
mengapa hasilnya menyimpang, daerah larik di-\eng{dump} pada
\reg{DS}:\hex{0000}:

\begin{session}[caption={\eng{Dump} larik \code{nilai} pada
\reg{DS}:\hex{0000}.}]
-d ds:0
13A6:0000  0A 14 1E 28 32 B7 00 00-00 00 00 00 00 00 00 00  ...(2...........
\end{session}

Lima byte pertama sesuai dengan data yang dideklarasikan, yaitu
\hex{0A}=\dec{10} sampai \hex{32}=\dec{50}. Byte keenam bernilai
\hex{B7}, yaitu sisa data lama di memori, bukan bagian dari larik.
Karena penghitung bernilai 6, byte itulah yang ikut dijumlahkan
sehingga akumulator menghasilkan \hex{014D} atau \dec{333} alih-alih
\hex{0096}. Inilah bukti yang mengubah dugaan menjadi kepastian:
kesalahannya bukan pada penjumlahan, melainkan pada batas perulangan.

\subsection{Membaca Bingkai Stack}

\eng{Stack} tumbuh ke arah alamat yang semakin kecil. Setiap
\asm{PUSH} dan setiap \asm{CALL} mengurangi \reg{SP} dan menuliskan
data pada \reg{SS}:\reg{SP} yang baru. Sebuah prosedur \eng{near}
meninggalkan susunan yang khas di dalam \eng{stack}, yang disebut
\term{bingkai stack} (\eng{stack frame}), seperti pada
Gambar~\ref{fig:stack-frame}.

\begin{figure}[H]
    \centering
    \begin{tikzpicture}[
        sel/.style={
            draw=obeBlue,
            line width=0.7pt,
            fill=obeBlueLight,
            minimum width=5.6cm,
            minimum height=0.72cm,
            font=\small,
            inner sep=2pt
        },
        selsp/.style={
            draw=obeMaroon,
            line width=1.0pt,
            fill=obeMaroonLight,
            minimum width=5.6cm,
            minimum height=0.72cm,
            font=\small,
            inner sep=2pt
        },
        al/.style={font=\scriptsize\ttfamily, text=obeGray, anchor=east},
        sp/.style={-{Stealth[length=2.5mm]}, line width=1.0pt, obeMaroon}
    ]
        \node[sel, fill=obeGrayLight] (d) at (0,0) {dasar \eng{stack}};
        \node[sel, below=0.10cm of d] (r) {alamat kembali \hex{0008}};
        \node[sel, below=0.10cm of r] (ax) {AX lama = \hex{13A6}};
        \node[sel, below=0.10cm of ax] (bx) {BX lama = \hex{0000}};
        \node[selsp, below=0.10cm of bx] (cx)
            {CX lama = \hex{00FF} \quad$\Leftarrow$ \reg{SP}};

        \node[al] at ([xshift=-0.18cm]d.west) {\hex{0100}};
        \node[al] at ([xshift=-0.18cm]r.west) {\hex{00FE}};
        \node[al] at ([xshift=-0.18cm]ax.west) {\hex{00FC}};
        \node[al] at ([xshift=-0.18cm]bx.west) {\hex{00FA}};
        \node[al] at ([xshift=-0.18cm]cx.west) {\hex{00F8}};

        \draw[sp] ([xshift=0.18cm]cx.east) -- ++(0.9cm,0)
            node[right, font=\scriptsize, text=obeMaroon, align=left]
            {\reg{SP} menunjuk\\puncak \eng{stack}};
    \end{tikzpicture}
    \caption{Bingkai \eng{stack} prosedur \asm{HITUNG} pada
    \file{bug-stack.asm} sesaat setelah ketiga \asm{PUSH}.}
    \label{fig:stack-frame}
\end{figure}

Alamat kembali berada \emph{di bawah} seluruh data yang ditekan
prosedur, sehingga ia hanya dapat diambil oleh \asm{RET} apabila
seluruh data di atasnya sudah diambil kembali. Bila jumlah \asm{PUSH}
dan \asm{POP} tidak seimbang, \reg{SP} berhenti pada posisi yang
salah. Sesi berikut memperlihatkan kejadian itu.

\begin{session}[caption={Menemukan \eng{stack} tidak seimbang melalui
\code{D} dan \code{T}.}, label={lst:sesi-stack}]
-g 0005
AX=13A6  BX=0000  CX=00FF  DX=0000  SP=0100  BP=0000  SI=0000  DI=0000
DS=13A6  ES=1391  SS=13A7  CS=13A1  IP=0005   NV UP EI PL NZ NA PO NC
13A1:0005 E80D00         CALL    0015

-t
AX=13A6  BX=0000  CX=00FF  DX=0000  SP=00FE  BP=0000  SI=0000  DI=0000
DS=13A6  ES=1391  SS=13A7  CS=13A1  IP=0015   NV UP EI PL NZ NA PO NC
13A1:0015 50             PUSH    AX

-t
AX=13A6  BX=0000  CX=00FF  DX=0000  SP=00FC  BP=0000  SI=0000  DI=0000
DS=13A6  ES=1391  SS=13A7  CS=13A1  IP=0016   NV UP EI PL NZ NA PO NC
13A1:0016 53             PUSH    BX

-t
AX=13A6  BX=0000  CX=00FF  DX=0000  SP=00FA  BP=0000  SI=0000  DI=0000
DS=13A6  ES=1391  SS=13A7  CS=13A1  IP=0017   NV UP EI PL NZ NA PO NC
13A1:0017 51             PUSH    CX

-t
AX=13A6  BX=0000  CX=00FF  DX=0000  SP=00F8  BP=0000  SI=0000  DI=0000
DS=13A6  ES=1391  SS=13A7  CS=13A1  IP=0018   NV UP EI PL NZ NA PO NC
13A1:0018 B80100         MOV     AX,0001

-d ss:00f8
13A7:00F8  FF 00 00 00 A6 13 08 00-00 00 00 00 00 00 00 00  ................

-t
AX=0001  BX=0000  CX=00FF  DX=0000  SP=00F8  BP=0000  SI=0000  DI=0000
DS=13A6  ES=1391  SS=13A7  CS=13A1  IP=001B   NV UP EI PL NZ NA PO NC
13A1:001B BB0200         MOV     BX,0002

-t
AX=0001  BX=0002  CX=00FF  DX=0000  SP=00F8  BP=0000  SI=0000  DI=0000
DS=13A6  ES=1391  SS=13A7  CS=13A1  IP=001E   NV UP EI PL NZ NA PO NC
13A1:001E 01D8           ADD     AX,BX

-t
AX=0003  BX=0002  CX=00FF  DX=0000  SP=00F8  BP=0000  SI=0000  DI=0000
DS=13A6  ES=1391  SS=13A7  CS=13A1  IP=0020   NV UP EI PL NZ NA PE NC
13A1:0020 89C1           MOV     CX,AX

-t
AX=0003  BX=0002  CX=0003  DX=0000  SP=00F8  BP=0000  SI=0000  DI=0000
DS=13A6  ES=1391  SS=13A7  CS=13A1  IP=0022   NV UP EI PL NZ NA PE NC
13A1:0022 59             POP     CX

-t
AX=0003  BX=0002  CX=00FF  DX=0000  SP=00FA  BP=0000  SI=0000  DI=0000
DS=13A6  ES=1391  SS=13A7  CS=13A1  IP=0023   NV UP EI PL NZ NA PE NC
13A1:0023 5B             POP     BX

-t
AX=0003  BX=0000  CX=00FF  DX=0000  SP=00FC  BP=0000  SI=0000  DI=0000
DS=13A6  ES=1391  SS=13A7  CS=13A1  IP=0024   NV UP EI PL NZ NA PE NC
13A1:0024 C3             RET

-d ss:00fc
13A7:00FC  A6 13 08 00-00 00 00 00 00 00 00 00 00 00 00 00  ................

-t
AX=0003  BX=0000  CX=00FF  DX=0000  SP=00FE  BP=0000  SI=0000  DI=0000
DS=13A6  ES=1391  SS=13A7  CS=13A1  IP=13A6   NV UP EI PL NZ NA PE NC
13A1:13A6 0000           ADD     [BX+SI],AL
\end{session}

\asm{CALL} menurunkan \reg{SP} dari \hex{0100} ke \hex{00FE} dan
menuliskan alamat kembali \hex{0008} di sana; ketiga \asm{PUSH}
kemudian menurunkan \reg{SP} sampai \hex{00F8}. \eng{Dump} pada
\reg{SS}:\hex{00F8} menunjukkan susunan yang persis sama dengan
Gambar~\ref{fig:stack-frame}. Setelah dua \asm{POP}, \reg{SP} berada
pada \hex{00FC} dan menunjuk \code{A6 13}, yaitu nilai \reg{AX} lama.
\asm{RET} mengambil kata itu sebagai alamat kembali, sehingga
\reg{IP} berubah menjadi \hex{13A6} --- jauh di luar batas program.
Baris terakhir sesi menunjukkan akibatnya: DEBUG membongkar isi memori
pada \hex{13A1}:\hex{13A6} dan menemukan instruksi sampah
\lstasm|ADD [BX+SI], AL|.

\subsection{Deteksi \eng{Overflow} melalui Flag dan \reg{SP}}

\eng{Overflow} aritmatika tidak selalu menghentikan program, tetapi
selalu meninggalkan jejak pada \eng{flag}: operasi yang melewati batas
bilangan bertanda menyalakan \flag{OF}, sedangkan yang melewati batas
bilangan tak bertanda menyalakan \flag{CF}. \eng{Overflow}
\eng{stack} meninggalkan jejak yang berbeda, yaitu nilai \reg{SP} yang
terus menurun tanpa pernah naik kembali. Bila perulangan menekan data
tanpa pernah mengambilnya, \eng{stack} akan melimpah dan menimpa
daerah lain; gejalanya adalah perubahan nilai variabel yang tidak
pernah disentuh program. Karena itu, kedua \eng{flag} tersebut dan
register \reg{SP} sebaiknya selalu disertakan pada tabel \eng{trace}
\cite{irvine1999}.

\subsection{Metodologi Debugging yang Sistematis}

\eng{Debugging} yang berhasil hampir selalu mengikuti alur yang sama.
Keenam tahapnya diperlihatkan pada
Gambar~\ref{fig:metodologi-debugging}.

\begin{figure}[H]
    \centering
    \begin{tikzpicture}[
        node distance=0.42cm,
        tahap/.style={
            draw=obeTeal,
            line width=0.8pt,
            fill=obeTealLight,
            rounded corners=3pt,
            text width=8.4cm,
            minimum height=1.0cm,
            align=center,
            inner sep=4pt,
            font=\small
        },
        panah/.style={-{Stealth[length=2.5mm]}, line width=0.9pt, obeTeal},
        balik/.style={-{Stealth[length=2.5mm]}, line width=0.9pt,
            obeMaroon, dashed}
    ]
        \node[tahap] (t1) {\textbf{1. Reproduksi}\\
            {\footnotesize\itshape pastikan kesalahan dapat dimunculkan ulang}};
        \node[tahap, below=of t1] (t2) {\textbf{2. Isolasi}\\
            {\footnotesize\itshape perkecil cakupan sampai bagian terkecil}};
        \node[tahap, below=of t2] (t3) {\textbf{3. Hipotesis}\\
            {\footnotesize\itshape rumuskan dugaan penyebab yang dapat diuji}};
        \node[tahap, below=of t3] (t4) {\textbf{4. Uji hipotesis}\\
            {\footnotesize\itshape buktikan dengan \code{T}, \code{G}, dan \code{D}}};
        \node[tahap, below=of t4] (t5) {\textbf{5. Perbaiki}\\
            {\footnotesize\itshape ubah satu hal saja pada satu waktu}};
        \node[tahap, below=of t5] (t6) {\textbf{6. Verifikasi}\\
            {\footnotesize\itshape ulangi pengujian pada seluruh kasus}};

        \draw[panah] (t1) -- (t2);
        \draw[panah] (t2) -- (t3);
        \draw[panah] (t3) -- (t4);
        \draw[panah] (t4) -- (t5);
        \draw[panah] (t5) -- (t6);

        \draw[balik] (t6.west) to[out=180, in=180, looseness=1.15]
            node[midway, left=2mm, font=\scriptsize, align=center,
                 text=obeMaroon]
            {hipotesis\\gugur atau\\perbaikan\\belum berhasil}
            (t3.west);
    \end{tikzpicture}
    \caption{Metodologi \eng{debugging} sistematis dan jalur pengulangan
    ketika hipotesis terbukti salah.}
    \label{fig:metodologi-debugging}
\end{figure}

Pada tahap isolasi, tentukan bagian program yang pasti bersalah dan
yang pasti benar. Pada tahap hipotesis, tuliskan dugaan dalam satu
kalimat yang dapat diuji, misalnya ``\reg{CX} tidak pernah mencapai
nol karena dinaikkan di dalam perulangan''. Pengujian dilakukan
dengan mengumpulkan bukti melalui tabel \eng{trace}. Pada tahap
perbaikan, ubah \emph{satu hal saja} agar penyebab keberhasilan dapat
dipastikan, lalu verifikasi dengan mengulang seluruh pengujian.

\begin{peringatanbox}
Kesalahan paling sering dalam \eng{debugging} adalah mengubah
beberapa baris sekaligus ketika program berperilaku salah. Bila
program kemudian berjalan benar, programmer tidak akan pernah tahu
baris mana yang sebenarnya memperbaiki keadaan, dan bug yang sama akan
terulang pada program berikutnya.
\end{peringatanbox}

\subsection{Studi Kasus Tiga Bug}

\begin{studikasusbox}
\textbf{Kasus 1: perulangan tak berujung karena \reg{CX} atau
\eng{flag} yang tidak berubah.}
Program \file{bug-loop.asm} pada Kode~\ref{lst:bug-loop} berhenti
tidak pada waktunya. Tabel~\ref{tab:trace-register} menjadi bukti
utamanya: pada langkah 6, \reg{CX} naik menjadi \hex{0006}, lalu
\asm{LOOP} menguranginya kembali menjadi \hex{0005}. Hipotesis
``\reg{CX} tidak pernah mencapai nol'' terbukti, dan perbaikannya
adalah menghapus \asm{INC CX} (Kode~\ref{lst:bug-loop-fix}).

Kerabat dekat kesalahan ini adalah perulangan berbasis \eng{flag}
yang menguji instruksi yang tidak mengubah \eng{flag}. Bila
\asm{SUB CX, 1} pada pola \asm{SUB}--\asm{JNZ} diganti dengan
\asm{MOV CX, BX}, maka \flag{ZF} tidak lagi mencerminkan penghitung
dan perulangan berjalan selamanya. Kaidah amannya: pastikan instruksi
yang mengubah \eng{flag} berada tepat sebelum instruksi yang menguji
\eng{flag} tersebut.
\end{studikasusbox}

\begin{studikasusbox}
\textbf{Kasus 2: kesalahan \eng{off-by-one} pada larik.}
Program pada Kode~\ref{lst:off-by-one} selesai dengan normal, sehingga
kesalahan sintaks dan kesalahan saat berjalan dapat disingkirkan.
\eng{Dump} \code{D DS:0} memperlihatkan bahwa byte keenam bernilai
\hex{B7}, sedangkan larik hanya berisi lima byte. Bukti itu cukup
untuk menyimpulkan bahwa penghitung perulangan melebihi jumlah elemen;
perbaikannya adalah mengubah \lstasm|MOV CX, 6| menjadi
\lstasm|MOV CX, 5|. Cara yang lebih aman adalah menghitung jumlah
elemen secara otomatis dengan \code{panjang EQU \$ - nilai} sehingga
kesalahan yang sama tidak terulang ketika larik diperpanjang.
\end{studikasusbox}

\begin{studikasusbox}
\textbf{Kasus 3: \eng{stack} tidak seimbang yang menyebabkan \asm{RET}
ke alamat salah.}
Program \file{bug-stack.asm} melompat ke alamat aneh. Gejala ini khas
\eng{stack} yang rusak, karena hanya \asm{CALL}, \asm{RET}, dan
\asm{INT} yang dapat memindahkan kendali ke luar batas program.
Langkah diagnostiknya adalah membandingkan jumlah \asm{PUSH} dengan
jumlah \asm{POP}, lalu membuktikannya melalui \eng{dump} \eng{stack}
seperti pada Kode~\ref{lst:sesi-stack}. Prosedur menekan tiga register
tetapi hanya memulihkan dua, sehingga \reg{SP} berhenti pada
\hex{00FC} dan \asm{RET} membaca \hex{13A6}. Perbaikan pada
Kode~\ref{lst:bug-stack-fix} menambahkan satu \asm{POP AX}.
\end{studikasusbox}

\begin{infobox}
Ketiga kasus di atas memiliki satu pola yang sama: gejalanya muncul di
satu tempat, sedangkan penyebabnya berada di tempat lain. Karena itu,
menelusuri program pada level instruksi --- bukan membaca ulang kode
sumber saja --- merupakan cara paling cepat menemukan penyebab
sesungguhnya \cite{ebrahim2018}.
\end{infobox}
